% Options for packages loaded elsewhere
\PassOptionsToPackage{unicode}{hyperref}
\PassOptionsToPackage{hyphens}{url}
\documentclass[
]{article}
\usepackage{xcolor}
\usepackage[margin=1in]{geometry}
\usepackage{amsmath,amssymb}
\setcounter{secnumdepth}{-\maxdimen} % remove section numbering
\usepackage{iftex}
\ifPDFTeX
  \usepackage[T1]{fontenc}
  \usepackage[utf8]{inputenc}
  \usepackage{textcomp} % provide euro and other symbols
\else % if luatex or xetex
  \usepackage{unicode-math} % this also loads fontspec
  \defaultfontfeatures{Scale=MatchLowercase}
  \defaultfontfeatures[\rmfamily]{Ligatures=TeX,Scale=1}
\fi
\usepackage{lmodern}
\ifPDFTeX\else
  % xetex/luatex font selection
\fi
% Use upquote if available, for straight quotes in verbatim environments
\IfFileExists{upquote.sty}{\usepackage{upquote}}{}
\IfFileExists{microtype.sty}{% use microtype if available
  \usepackage[]{microtype}
  \UseMicrotypeSet[protrusion]{basicmath} % disable protrusion for tt fonts
}{}
\makeatletter
\@ifundefined{KOMAClassName}{% if non-KOMA class
  \IfFileExists{parskip.sty}{%
    \usepackage{parskip}
  }{% else
    \setlength{\parindent}{0pt}
    \setlength{\parskip}{6pt plus 2pt minus 1pt}}
}{% if KOMA class
  \KOMAoptions{parskip=half}}
\makeatother
\usepackage{color}
\usepackage{fancyvrb}
\newcommand{\VerbBar}{|}
\newcommand{\VERB}{\Verb[commandchars=\\\{\}]}
\DefineVerbatimEnvironment{Highlighting}{Verbatim}{commandchars=\\\{\}}
% Add ',fontsize=\small' for more characters per line
\usepackage{framed}
\definecolor{shadecolor}{RGB}{248,248,248}
\newenvironment{Shaded}{\begin{snugshade}}{\end{snugshade}}
\newcommand{\AlertTok}[1]{\textcolor[rgb]{0.94,0.16,0.16}{#1}}
\newcommand{\AnnotationTok}[1]{\textcolor[rgb]{0.56,0.35,0.01}{\textbf{\textit{#1}}}}
\newcommand{\AttributeTok}[1]{\textcolor[rgb]{0.13,0.29,0.53}{#1}}
\newcommand{\BaseNTok}[1]{\textcolor[rgb]{0.00,0.00,0.81}{#1}}
\newcommand{\BuiltInTok}[1]{#1}
\newcommand{\CharTok}[1]{\textcolor[rgb]{0.31,0.60,0.02}{#1}}
\newcommand{\CommentTok}[1]{\textcolor[rgb]{0.56,0.35,0.01}{\textit{#1}}}
\newcommand{\CommentVarTok}[1]{\textcolor[rgb]{0.56,0.35,0.01}{\textbf{\textit{#1}}}}
\newcommand{\ConstantTok}[1]{\textcolor[rgb]{0.56,0.35,0.01}{#1}}
\newcommand{\ControlFlowTok}[1]{\textcolor[rgb]{0.13,0.29,0.53}{\textbf{#1}}}
\newcommand{\DataTypeTok}[1]{\textcolor[rgb]{0.13,0.29,0.53}{#1}}
\newcommand{\DecValTok}[1]{\textcolor[rgb]{0.00,0.00,0.81}{#1}}
\newcommand{\DocumentationTok}[1]{\textcolor[rgb]{0.56,0.35,0.01}{\textbf{\textit{#1}}}}
\newcommand{\ErrorTok}[1]{\textcolor[rgb]{0.64,0.00,0.00}{\textbf{#1}}}
\newcommand{\ExtensionTok}[1]{#1}
\newcommand{\FloatTok}[1]{\textcolor[rgb]{0.00,0.00,0.81}{#1}}
\newcommand{\FunctionTok}[1]{\textcolor[rgb]{0.13,0.29,0.53}{\textbf{#1}}}
\newcommand{\ImportTok}[1]{#1}
\newcommand{\InformationTok}[1]{\textcolor[rgb]{0.56,0.35,0.01}{\textbf{\textit{#1}}}}
\newcommand{\KeywordTok}[1]{\textcolor[rgb]{0.13,0.29,0.53}{\textbf{#1}}}
\newcommand{\NormalTok}[1]{#1}
\newcommand{\OperatorTok}[1]{\textcolor[rgb]{0.81,0.36,0.00}{\textbf{#1}}}
\newcommand{\OtherTok}[1]{\textcolor[rgb]{0.56,0.35,0.01}{#1}}
\newcommand{\PreprocessorTok}[1]{\textcolor[rgb]{0.56,0.35,0.01}{\textit{#1}}}
\newcommand{\RegionMarkerTok}[1]{#1}
\newcommand{\SpecialCharTok}[1]{\textcolor[rgb]{0.81,0.36,0.00}{\textbf{#1}}}
\newcommand{\SpecialStringTok}[1]{\textcolor[rgb]{0.31,0.60,0.02}{#1}}
\newcommand{\StringTok}[1]{\textcolor[rgb]{0.31,0.60,0.02}{#1}}
\newcommand{\VariableTok}[1]{\textcolor[rgb]{0.00,0.00,0.00}{#1}}
\newcommand{\VerbatimStringTok}[1]{\textcolor[rgb]{0.31,0.60,0.02}{#1}}
\newcommand{\WarningTok}[1]{\textcolor[rgb]{0.56,0.35,0.01}{\textbf{\textit{#1}}}}
\usepackage{graphicx}
\makeatletter
\newsavebox\pandoc@box
\newcommand*\pandocbounded[1]{% scales image to fit in text height/width
  \sbox\pandoc@box{#1}%
  \Gscale@div\@tempa{\textheight}{\dimexpr\ht\pandoc@box+\dp\pandoc@box\relax}%
  \Gscale@div\@tempb{\linewidth}{\wd\pandoc@box}%
  \ifdim\@tempb\p@<\@tempa\p@\let\@tempa\@tempb\fi% select the smaller of both
  \ifdim\@tempa\p@<\p@\scalebox{\@tempa}{\usebox\pandoc@box}%
  \else\usebox{\pandoc@box}%
  \fi%
}
% Set default figure placement to htbp
\def\fps@figure{htbp}
\makeatother
\setlength{\emergencystretch}{3em} % prevent overfull lines
\providecommand{\tightlist}{%
  \setlength{\itemsep}{0pt}\setlength{\parskip}{0pt}}
\usepackage{bookmark}
\IfFileExists{xurl.sty}{\usepackage{xurl}}{} % add URL line breaks if available
\urlstyle{same}
\hypersetup{
  pdftitle={The Probability Density Function},
  hidelinks,
  pdfcreator={LaTeX via pandoc}}

\title{The Probability Density Function}
\author{}
\date{\vspace{-2.5em}2026-04-17}

\begin{document}
\maketitle

\subsection{Introduction}\label{introduction}

For a continuous random variable, probability is distributed smoothly
across the support rather than concentrated on countable points. The
\textbf{probability density function (PDF)} encodes this distribution:
integrating the PDF over a set gives the probability that the variable
lands in that set.

\subsection{Prerequisites}\label{prerequisites}

Calculus (integration) and the concept of a continuous random variable.

\subsection{Theory}\label{theory}

A continuous random variable \(X\) has a PDF \(f_X\) if, for every
measurable set \(A\),

\[P(X \in A) = \int_A f_X(x) \, dx.\]

A valid PDF satisfies:

\begin{enumerate}
\def\labelenumi{\arabic{enumi}.}
\tightlist
\item
  \(f_X(x) \geq 0\) for all \(x\).
\item
  \(\int_{-\infty}^{\infty} f_X(x) \, dx = 1\).
\end{enumerate}

Crucially, \(f_X(x)\) is \textbf{not} a probability. For continuous
\(X\), \(P(X = x) = 0\) for every \(x\). The density can exceed 1 at a
given point, provided the integral is bounded. Density has units of
inverse \(x\): for \(x\) in seconds, \(f_X(x)\) has units of
\(1/\text{second}\).

Useful relationships:

\begin{itemize}
\tightlist
\item
  CDF: \(F_X(x) = \int_{-\infty}^x f_X(u) \, du\).
\item
  PDF from CDF: \(f_X(x) = F_X'(x)\) where \(F_X\) is differentiable.
\item
  Expected value: \(E[X] = \int x f_X(x) \, dx\).
\item
  Variance: \(\mathrm{Var}(X) = \int (x - E[X])^2 f_X(x) \, dx\).
\end{itemize}

In R, each standard continuous distribution has a \texttt{d*()} function
returning its PDF.

\subsection{Assumptions}\label{assumptions}

The variable is continuous (absolutely continuous with respect to
Lebesgue measure). Mixed or discrete variables do not have ordinary PDFs
(though the mixed case has a generalised density with respect to a mixed
measure).

\subsection{R Implementation}\label{r-implementation}

\begin{Shaded}
\begin{Highlighting}[]
\FunctionTok{library}\NormalTok{(ggplot2)}

\NormalTok{x\_grid }\OtherTok{\textless{}{-}} \FunctionTok{seq}\NormalTok{(}\SpecialCharTok{{-}}\DecValTok{4}\NormalTok{, }\DecValTok{4}\NormalTok{, }\AttributeTok{length.out =} \DecValTok{801}\NormalTok{)}
\NormalTok{fx }\OtherTok{\textless{}{-}} \FunctionTok{dnorm}\NormalTok{(x\_grid, }\AttributeTok{mean =} \DecValTok{0}\NormalTok{, }\AttributeTok{sd =} \DecValTok{1}\NormalTok{)}

\FunctionTok{ggplot}\NormalTok{(}\FunctionTok{data.frame}\NormalTok{(}\AttributeTok{x =}\NormalTok{ x\_grid, }\AttributeTok{f =}\NormalTok{ fx), }\FunctionTok{aes}\NormalTok{(x, f)) }\SpecialCharTok{+}
  \FunctionTok{geom\_line}\NormalTok{(}\AttributeTok{colour =} \StringTok{"\#2A9D8F"}\NormalTok{, }\AttributeTok{linewidth =} \DecValTok{1}\NormalTok{) }\SpecialCharTok{+}
  \FunctionTok{geom\_area}\NormalTok{(}\AttributeTok{data =} \FunctionTok{subset}\NormalTok{(}\FunctionTok{data.frame}\NormalTok{(}\AttributeTok{x =}\NormalTok{ x\_grid, }\AttributeTok{f =}\NormalTok{ fx),}
\NormalTok{                          x }\SpecialCharTok{\textgreater{}=} \SpecialCharTok{{-}}\FloatTok{1.96} \SpecialCharTok{\&}\NormalTok{ x }\SpecialCharTok{\textless{}=} \FloatTok{1.96}\NormalTok{),}
            \FunctionTok{aes}\NormalTok{(}\AttributeTok{y =}\NormalTok{ f), }\AttributeTok{fill =} \StringTok{"\#2A9D8F"}\NormalTok{, }\AttributeTok{alpha =} \FloatTok{0.25}\NormalTok{) }\SpecialCharTok{+}
  \FunctionTok{labs}\NormalTok{(}\AttributeTok{x =} \StringTok{"x"}\NormalTok{, }\AttributeTok{y =} \StringTok{"f(x)"}\NormalTok{,}
       \AttributeTok{title =} \StringTok{"Standard normal PDF with central 95\% shaded"}\NormalTok{) }\SpecialCharTok{+}
  \FunctionTok{theme\_minimal}\NormalTok{()}

\CommentTok{\# Verify normalisation by numerical integration}
\FunctionTok{integrate}\NormalTok{(dnorm, }\SpecialCharTok{{-}}\ConstantTok{Inf}\NormalTok{, }\ConstantTok{Inf}\NormalTok{)}

\CommentTok{\# Probability of an interval via integration}
\FunctionTok{integrate}\NormalTok{(dnorm, }\SpecialCharTok{{-}}\FloatTok{1.96}\NormalTok{, }\FloatTok{1.96}\NormalTok{)        }\CommentTok{\# about 0.95}

\CommentTok{\# Equivalently via the CDF}
\FunctionTok{pnorm}\NormalTok{(}\FloatTok{1.96}\NormalTok{) }\SpecialCharTok{{-}} \FunctionTok{pnorm}\NormalTok{(}\SpecialCharTok{{-}}\FloatTok{1.96}\NormalTok{)}

\CommentTok{\# Moments via integration}
\FunctionTok{integrate}\NormalTok{(}\ControlFlowTok{function}\NormalTok{(x) x }\SpecialCharTok{*} \FunctionTok{dnorm}\NormalTok{(x), }\SpecialCharTok{{-}}\ConstantTok{Inf}\NormalTok{, }\ConstantTok{Inf}\NormalTok{)        }\CommentTok{\# E[X]}
\FunctionTok{integrate}\NormalTok{(}\ControlFlowTok{function}\NormalTok{(x) x}\SpecialCharTok{\^{}}\DecValTok{2} \SpecialCharTok{*} \FunctionTok{dnorm}\NormalTok{(x), }\SpecialCharTok{{-}}\ConstantTok{Inf}\NormalTok{, }\ConstantTok{Inf}\NormalTok{)      }\CommentTok{\# E[X\^{}2]}
\end{Highlighting}
\end{Shaded}

\subsection{Output \& Results}\label{output-results}

\begin{verbatim}
1 with absolute error < 9.4e-05

0.9500042 with absolute error < 1e-11

[1] 0.9500042

-4.4e-11 with absolute error < ...   # E[X] ≈ 0
1 with absolute error < ...          # E[X^2] = Var(X) for standard normal
\end{verbatim}

Integrating the standard normal PDF over \([-1.96, 1.96]\) gives 0.95 --
the well-known central probability -- and the \texttt{pnorm} difference
recovers the same value via the CDF.

\subsection{Interpretation}\label{interpretation}

For applied purposes, a PDF is visualised by its curve; probabilities
are computed via \texttt{p*()} functions in R rather than by explicit
integration. When reporting a continuous distribution in a paper, the
density plot, mean, SD, and median are the standard summaries.

\subsection{Practical Tips}\label{practical-tips}

\begin{itemize}
\tightlist
\item
  Density at a point is not probability; do not compare \(f(x_1)\) to
  \(f(x_2)\) as probabilities of outcomes 1 and 2. Compare
  \(P(X \in (x_1 - h, x_1 + h))\) for small \(h\) instead.
\item
  The density can be unbounded (e.g., Gamma with shape \(< 1\) near
  \(0\)); integrability, not boundedness, is what matters.
\item
  Kernel density estimates (\texttt{density()}) approximate the PDF from
  a sample.
\item
  Log-densities (\texttt{dnorm(...,\ log\ =\ TRUE)}) avoid underflow
  when evaluating likelihoods with many observations.
\item
  The relationship \(f = F'\) is the starting point for the inverse-CDF
  method: to generate \(X\), apply \(F^{-1}\) to a uniform random
  variable.
\end{itemize}

\subsection{For Reviewers}\label{for-reviewers}

\textbf{What to look for in a paper using this method.}

\begin{itemize}
\tightlist
\item
  \textbf{Common misapplications.}
\item
  \textbf{Diagnostics that should be reported but often aren't.}
\item
  \textbf{Red flags in tables and figures.}
\item
  \textbf{What to verify.}
\item
  \textbf{What an adequate Methods paragraph must contain.}
\end{itemize}

\subsection{See also --- labs in this
chapter}\label{see-also-labs-in-this-chapter}

\begin{itemize}
\tightlist
\item
  \href{labs/lab-probability-conditional-probability-bayes-theorem.Rmd}{Probability,
  conditional probability, Bayes' theorem}
\item
  \href{labs/lab-discrete-distributions-bernoulli-binomial-poisson.Rmd}{Discrete
  distributions: Bernoulli, binomial, Poisson}
\item
  \href{labs/lab-continuous-distributions-and-q-q-plots.Rmd}{Continuous
  distributions and Q-Q plots}
\end{itemize}

\end{document}
